% !TeX program = pdflatex
\documentclass[12pt,letterpaper]{article}
\usepackage[utf8]{inputenc}
\usepackage[letterpaper, margin=1in, bottom=2.5in, top=2.2in]{geometry}
\usepackage{fancyhdr}
\setlength{\headheight}{15pt}
\setlength{\footskip}{72pt}

% Import all the same packages and setup as Parts 1-2
\usepackage{xcolor}
\usepackage{amsmath}
\usepackage{anyfontsize}
\usepackage{graphicx}
\usepackage{tikz}
\usepackage{listings}
\usepackage{sourcesanspro}
\renewcommand{\familydefault}{\sfdefault}
\usepackage{fontawesome5}
\usepackage{titlesec}
\usepackage{setspace}
\usepackage{hyperref}
\usepackage{mdframed}
\usepackage{qrcode}
\usetikzlibrary{calc,shapes,positioning}
\usepackage{eso-pic}

% Command for inline code highlighting
\newcommand{\code}[1]{\texttt{\textcolor{accentColor}{#1}}}
\newcommand{\codebf}[1]{\texttt{\textcolor{codeKeywordColor}{#1}}}

\AtBeginDocument{\color{primaryColor}}

% Colors updated with JavaScript brand colors
\definecolor{bgColor}{RGB}{13, 17, 23}  % Dark background similar to JS docs
\definecolor{primaryColor}{RGB}{255, 255, 255}  % White text for better contrast
\definecolor{accentColor}{RGB}{255, 220, 60}  % Brighter JavaScript Yellow
\definecolor{jsSecondary}{RGB}{180, 180, 180}  % Light gray as secondary color
\definecolor{secondaryColor}{RGB}{230, 235, 240}  % Almost white gray
\definecolor{terminalBg}{RGB}{22, 22, 22}  % Terminal background
\definecolor{terminalFrame}{RGB}{40, 40, 40}  % Terminal frame
\definecolor{lineNumberColor}{RGB}{100, 100, 100}  % Line number color
\definecolor{dividerColor}{RGB}{120, 130, 150}  % Brighter divider color

% Code colors
\definecolor{codeTextColor}{RGB}{255, 255, 255}  % Pure white for basic text
\definecolor{codeKeywordColor}{RGB}{97, 175, 239}  % Light blue for keywords
\definecolor{codeCommentColor}{RGB}{121, 192, 120}  % Green for comments
\definecolor{codeStringColor}{RGB}{255, 220, 60}  % JS Yellow for strings
\definecolor{codeMethodColor}{RGB}{255, 160, 122}  % Light orange for methods
\definecolor{codeFunctionColor}{RGB}{230, 180, 80}  % Golden yellow for functions
\definecolor{codeNumberColor}{RGB}{180, 230, 180}  % Light green for numbers


\pagecolor{bgColor}

\hypersetup{
    colorlinks=true,
    linkcolor=accentColor,
    filecolor=accentColor,
    urlcolor=accentColor,
}

% Define JavaScript language for listings with async/await keywords
\lstdefinelanguage{JavaScript}{
  keywords={break, case, catch, continue, debugger, default, delete, do, else, false, finally, for, function, if, in, instanceof, new, null, return, switch, this, throw, true, try, typeof, var, void, while, with, let, const, class, export, import, yield, async, await, Promise, resolve, reject},
  sensitive=true,
  comment=[l]{//},
  morecomment=[s]{/*}{*/},
  morestring=[b]',
  morestring=[b]"
}

% Listings configuration using JavaScript as default
\lstset{
  language=JavaScript,
  basicstyle=\ttfamily\bfseries\color{codeTextColor},
  backgroundcolor=\color{terminalBg},
  commentstyle=\color{codeCommentColor},
  keywordstyle=\color{codeKeywordColor},
  stringstyle=\color{codeStringColor},
  numberstyle=\color{lineNumberColor},
  breaklines=true,
  breakatwhitespace=true,
  tabsize=2,
  showstringspaces=false,
  frame=none,
  xleftmargin=15pt,
  xrightmargin=0pt,
  aboveskip=5pt,
  belowskip=5pt,
  numbers=left,
  numbersep=8pt,
  extendedchars=true,
  keepspaces=true,
  columns=flexible,
  lineskip=2pt,
  emph={[2]require,createReadStream,createWriteStream,pipe,on,once,createGzip,pipeline,Transform,Readable,Writable,Duplex, setTimeout, console, Error, Math, JSON},
  emphstyle={[2]\color{codeFunctionColor}}
}

\newenvironment{macterminal}{%
    \begin{mdframed}[
        linecolor=terminalFrame,
        backgroundcolor=terminalBg,
        roundcorner=5pt,
        skipabove=5pt,
        skipbelow=5pt,
        linewidth=1pt,
        innertopmargin=5pt,
        frametitle={%
            \tikz[baseline=(current bounding box.east), outer sep=0pt]{
                \fill[red!80!black] (0,0) circle (5pt);
                \fill[yellow!80!black] (0.7,0) circle (5pt);
                \fill[green!70!black] (1.4,0) circle (5pt);
            }
        },
        frametitlealignment=\raggedright,
        frametitleaboveskip=8pt,
        frametitlebelowskip=0pt,
    ]
}{%
    \end{mdframed}%
}

\newcommand{\verspace}{\vspace{5pt}}
\newcommand{\languagetag}[1]{
    \begin{tikzpicture}[baseline]
        \node[fill=accentColor, text=bgColor, rounded corners=5pt, inner sep=7pt] {
            {\normalsize\textbf{#1}}
        };
    \end{tikzpicture}
}

% Title formatting
\titleformat{\section}
  {\LARGE\bfseries\color{primaryColor}}
  {\thesection. }
  {0pt}
  {}
  []

\titleformat{\subsection}
  {\Large\bfseries\color{accentColor}}
  {\thesubsection. }
  {0pt}
  {}
  []

\titleformat{\subsubsection}
  {\large\bfseries\color{primaryColor}}
  {\thesubsubsection. }
  {0pt}
  {}
  []

\titlespacing*{\section}{0pt}{18pt}{8pt}
\titlespacing*{\subsection}{0pt}{12pt}{5pt}
\titlespacing*{\subsubsection}{0pt}{8pt}{3pt}

% Header and footer setup
\pagestyle{fancy}
\fancyhf{}

\fancyhead[L]{
    \begin{tikzpicture}[remember picture, overlay]
        \fill[accentColor, rounded corners=3pt] (0,0) rectangle (2.2cm,0.7cm);
        \node[text=bgColor, font=\bfseries] at (1.1cm,0.35cm) {NODE.JS};
        \fill[secondaryColor] (2.35cm,0.1cm) rectangle (2.40cm,0.6cm);
        \node[text=primaryColor, font=\bfseries, anchor=west] at (2.35cm,0.35cm) {PROMISES};
    \end{tikzpicture}
}

\renewcommand{\headrulewidth}{0pt}
\renewcommand{\footrulewidth}{0.2pt}
\renewcommand{\footruleskip}{1cm}

\fancyfoot[C]{
    \vspace*{0.1cm}
    \noindent
    \begin{minipage}{\textwidth}
        \begin{flushleft}
            % Profile photo
            \raisebox{0.7cm}{
            \begin{tikzpicture}[baseline]
                \path[fill=bgColor] (0,0) circle (0.6cm);
                \clip (0,0) circle (0.6cm);
                \node at (0,0) {
                    \includegraphics[width=1.2cm,height=1.2cm]{../../Images/profile-image.jpeg}
                };
            \end{tikzpicture}
            }
            % Profile information
            \begin{minipage}[b]{0.5\textwidth}
                {\large\bfseries\color{primaryColor}Alejandro Sánchez Yalí}
                \par\vspace{1pt}
                {\small\color{secondaryColor}Software Developer | AI \& Blockchain Enthusiast}
                \par\vspace{1pt}
                {\small\color{accentColor}\faGlobe\hspace{5pt}\color{secondaryColor}www.asanchezyali.com}
            \end{minipage}
        \end{flushleft}
        \vspace{6pt}
    \end{minipage}
}

\fancypagestyle{plain}{
    \fancyhf{}
    \renewcommand{\headrulewidth}{0pt}
    \renewcommand{\footrulewidth}{0pt}
}

\renewcommand{\labelitemi}{\textcolor{accentColor}{$\bullet$}}
\renewcommand{\labelitemii}{\textcolor{secondaryColor}{$\circ$}}

\usepackage{relsize}
\AtBeginDocument{\relsize{1}}

% Command for elegant QR
\newcommand{\elegantqr}[2]{
    \qrcode[height=2.5cm]{#1}
    \\[0.1cm]
    {\hspace{0.2cm}\color{primaryColor}\small #2\par}
}

\newcommand{\BackgroundPic}{%
    \put(0,0){%
        \begin{tikzpicture}[remember picture, overlay]
            % Place the background image covering the entire page
            \node[inner sep=0pt] at (current page.center) {
                \includegraphics[width=\paperwidth,height=\paperheight,keepaspectratio=false]{../../Images/Intense Focus at Gaming Event.jpeg}
            };
            
            % Add a semi-transparent dark layer over the image
            % You can adjust the color and opacity as needed
            \fill[black, opacity=0.7] (current page.south west) rectangle (current page.north east);
            
            % If you want a dark gradient instead of a solid color, use this:
            % \shade[top color=black!80, bottom color=black!40, opacity=0.7] 
            %    (current page.south west) rectangle (current page.north east);
        \end{tikzpicture}%
    }%
}

\newcommand{\titlepagecontents}{%
    \AddToShipoutPicture*{\BackgroundPic}
    \vspace*{2cm}
    \begin{flushleft}
    \languagetag{JavaScript}\\[0.4cm] % Use JS tag
    {\fontsize{38}{52}\bfseries\color{primaryColor}Understanding \color{accentColor}JavaScript\\\color{accentColor}Promises\par}
    \vspace{0.3cm}
    {\fontsize{18}{52}\color{secondaryColor}Handling Asynchronous Operations Like a Pro \\ Part 1/3\par}
    \vspace{0.3cm}
    {\color{secondaryColor}\today\par}
    \vspace{1.8cm}
    \elegantqr{https://github.com/asanchezyali/social-media-posts/tree/main/JavaScript/Promises}{Source Code}
    \end{flushleft}
}

\newcommand{\finalpagecontents}{%
    \AddToShipoutPicture*{\BackgroundPic}
    \vspace*{3cm}
    % Left-aligned elements with EXACT same structure as title page
    \begin{flushleft}
    \languagetag{Feedback}\\[0.4cm]
    {\fontsize{46}{52}\bfseries\color{primaryColor}Found this \color{accentColor}helpful?\par}
    \vspace{0.3cm}
    {\fontsize{18}{52}\color{secondaryColor}Save, comment and share\par}  
    \vspace{0.3cm}
    {\color{secondaryColor}\today\par}
    \end{flushleft}
}

\begin{document}
\setlength{\parindent}{0pt}
\begin{titlepage}
    \titlepagecontents
\end{titlepage}

\section{What's the Deal with Promises?}

Asynchronous operations are fundamental in modern JavaScript, enabling tasks like fetching data, reading files, or handling user interactions without blocking the main thread. However, managing the flow and potential errors of these operations can be complex. This is where \code{Promises} provide a robust solution.

\subsection{Why Promises?}

Consider an operation that takes time to complete, such as requesting data from an API. A synchronous approach would halt script execution until the data arrives, leading to unresponsive applications. Promises offer a way to handle such asynchronous tasks effectively.

They act as placeholders for a future value, representing the eventual result of an asynchronous operation. Instead of blocking, the operation is initiated, and the promise object is returned immediately. This allows the rest of the script to continue running.

Promises improve upon older asynchronous patterns like callbacks by providing:

\begin{itemize}
    \item A clearer, more manageable structure for handling results or errors.
    \item Better composability for chaining multiple asynchronous operations.
    \item A standardized way to manage asynchronous flow, reducing complexity (often referred to as "callback hell").
\end{itemize}

A Promise exists in one of three states:
\begin{itemize}
    \item \codebf{Pending:} The initial state; the asynchronous operation has not yet completed.
    \item \codebf{Fulfilled (Resolved):} The operation completed successfully, and the promise holds the resulting value.
    \item \codebf{Rejected:} The operation failed, and the promise holds the reason (typically an error object).
\end{itemize}

Once a promise transitions from pending to either fulfilled or rejected, it becomes \textbf{settled}, and its state and value (or reason) become immutable.

\section{Creating Your First Promise}

Let's make a simple promise. We'll simulate an async task using \code{setTimeout}.

\textbf{How to Run:}
\begin{itemize}
    \item Save the code below as \verb|01_creating_promise.js|.
    \item Open your terminal and run: \verb|node 01_creating_promise.js|
    \item Or, paste the code directly into your browser's developer console.
\end{itemize}

\begin{macterminal}
  \lstinputlisting[language=JavaScript]{codes/01_creating_promise.js}
\end{macterminal}

When you run this:
\begin{itemize}
    \item The \code{new Promise(...)} creates the promise.
    \item The executor function \code{(resolve, reject) => \{\...\}} runs immediately.
    \item The \code{setTimeout} schedules the success/failure logic to run after 2 seconds.
    \item The code continues, logging "Promise created...". The promise is \textbf{pending}.
    \item After 2 seconds, \code{setTimeout}'s callback runs, calling either \code{resolve} (making the promise fulfilled) or \code{reject} (making it rejected).
\end{itemize}

But how do we actually *use* the result or handle the error?

\section{Handling Promises: \code{.then()}, \code{.catch()}, \code{.finally()}}

Okay, we made a promise. Now, how do we react when it settles (fulfills or rejects)? We use special methods attached to the promise:

\begin{itemize}
    \item \codebf{.then(onFulfilled, onRejected)}: Attaches callbacks for when the promise is fulfilled (first argument) or rejected (second argument - less common, usually use \code{.catch()}).
    \item \codebf{.catch(onRejected)}: Attaches a callback specifically for when the promise is rejected. It's like a \code{try...catch} block for promises.
    \item \codebf{.finally(onFinally)}: Attaches a callback that runs \textit{always}, whether the promise was fulfilled or rejected. Great for cleanup tasks (like hiding a loading spinner).
\end{itemize}

Let's handle the promise we created earlier:

\textbf{How to Run:}
\begin{itemize}
    \item Save the code below as \verb|02_then_catch_finally.js|.
    \item Open your terminal and run: \verb|node 02_then_catch_finally.js|
    \item Or, paste the code directly into your browser's developer console.
\end{itemize}

\begin{macterminal}
  \lstinputlisting[language=JavaScript]{codes/02_then_catch_finally.js}
\end{macterminal}

Key takeaways:
\begin{itemize}
    \item Handlers (\code{.then}, \code{.catch}, \code{.finally}) are attached to the promise object to react to its settlement.
    \item These handlers execute asynchronously when the promise settles, not immediately upon attachment.
    \item \code{.then()} receives the resolved value as its argument.
    \item \code{.catch()} receives the rejection reason (usually an Error object) as its argument.
    \item \code{.finally()} receives no arguments and executes regardless of the outcome.
\end{itemize}

This forms the foundation of promise-based asynchronous programming. Their true power becomes more apparent when chaining multiple operations, which will be covered subsequently.

\section{Conclusions about JavaScript Promises}

In this first part of our exploration of Promises in JavaScript, we have learned:

\begin{itemize}
    \item \textbf{\textcolor{accentColor}{Fundamentals:}} Promises are objects that represent the eventual result of an asynchronous operation, allowing for cleaner and more structured code than traditional callbacks.
    
    \item \textbf{\textcolor{accentColor}{States:}} A promise can be in one of three states: pending, fulfilled, or rejected. Once a promise is settled (fulfilled or rejected), its state and value become immutable.
    
    \item \textbf{\textcolor{accentColor}{Result Handling:}} The \code{.then()}, \code{.catch()} and \code{.finally()} methods provide a clear interface for handling both successful results and errors from asynchronous operations.
    
    \item \textbf{\textcolor{accentColor}{Organized Asynchrony:}} Promises keep code readable even when working with operations that don't complete immediately, making it easier to manage asynchronous execution flow.
\end{itemize}

Promises are fundamental to modern JavaScript development and form the foundation for more advanced features like \code{async/await}. In part two, we'll explore promise chaining, methods like \code{Promise.all()} and \code{Promise.race()}, and advanced patterns that unlock the full potential of asynchronous programming in JavaScript.

\section{References}

\begin{itemize}
    \item MDN Web Docs. (2025). \textit{Using promises}. \href{https://developer.mozilla.org/en-US/docs/Web/JavaScript/Guide/Using_promises}{Link}
    
    \item MDN Web Docs. (2025). \textit{Promise}. \href{https://developer.mozilla.org/en-US/docs/Web/JavaScript/Reference/Global_Objects/Promise}{Link}
    
    \item ECMA International. (2026). \textit{ECMAScript 2026 Language Specification: Promise Objects}. \href{https://tc39.es/ecma262/multipage/control-abstraction-objects.html#sec-promise-objects}{Link}
    
    \item javascript.info. (2025). \textit{Promise}. \href{https://javascript.info/promise-basics}{Link}
    
    \item Parker, D. (2015). \textit{JavaScript with Promises: Managing Asynchronous Code}. O'Reilly Media.

    \item This article was translated, edited and written in collaboration with AI. If you find any inconsistencies or have suggestions for improvement, please don't hesitate to open an issue in our GitHub repository at \href{https://github.com/asanchezyali/social-media-posts}{github} or reach out directly.
\end{itemize}

\section{Explore My Other Posts}
\vspace{10pt}
\begin{tikzpicture}
  % Section background - increased height slightly to ensure all content fits
  \fill[color=terminalBg, rounded corners=5pt] (0,0) rectangle (\textwidth,-5);
  
  % Calculate width available for text (total width minus QR code width minus margins)
  \pgfmathsetmacro{\availableWidth}{\textwidth-4cm}
  
  % Left column content
  \begin{scope}
    % Section title
    \node[text=accentColor, font=\Large\bfseries, anchor=north west] at (0.5cm,-0.5cm) {Enjoyed This Content?};
    
    % Subtitle
    \node[text=primaryColor, font=\normalsize, anchor=north west] at (0.5cm,-1.2cm) {Don't miss my previous post about:};
    
    % Node.js Streams Part 1 title - with highlighted style and limited width
    \node[text=secondaryColor, font=\large\bfseries, anchor=north west, text width=\availableWidth] at (0.5cm,-1.9cm) {\color{accentColor}Node.js Streams: \color{primaryColor}Part 2: Types \& Advanced Operations};
    
    % Brief description of the post - with limited width and adjusted position
    \node[text=jsSecondary, font=\normalsize, text width=\availableWidth, anchor=north west] at (0.5cm,-3.1cm) {
      Learn the fundamentals of Node.js Streams and discover how they can dramatically reduce memory usage when processing large files.
    };
  \end{scope}
  
  % QR Code on the right - vertically centered
  \node[anchor=center] at ({\textwidth-1.5cm}, {-2.5cm}) {
    \qrcode[height=2.2cm]{https://www.linkedin.com/feed/update/urn:li:activity:7312444909154623489/}
  };
  
\end{tikzpicture}
\vspace{5pt}

\clearpage
\thispagestyle{empty}
\finalpagecontents

\end{document}
